\chapter{Loop-Transformation Kernels}
\label{ch:loop-kernels}
% Data: llr40Matrix/ (HPCAgent-Bench tag llr-focus40; v2 measured at benchmark commit 26a4f0cf).

The second kernel set focuses on loop transformations. It contains the 40 kernels of
HPCAgent-Bench that are built to test loop transformations, such as interchanging, splitting or
fusing loops. We call it the loop-transformation set, or LLR-40 after the benchmark's tag
for it.

\section{Kernel Set}
\label{sec:loop-kernel-set}

Twenty-nine of the kernels come from TSVC-2~\cite{tsvc2}, the C version of the classic test
suite for vectorising compilers~\cite{callahan1988tsvc,maleki2011evaluation}; three of these
carry new names. The other eleven were written for the benchmark and add patterns that TSVC
lacks, such as fusion through a temporary array, scans, scatters with duplicate indices, and
wavefronts.

Each kernel is written once in NumPy, as explicit loops; this version defines the correct
result. The benchmark's translators turn it into C, C++, Fortran and Numba code. For 35 kernels
there is also a hand-written C version, ported from TSVC-2. All kernels compute in
\texttt{float64}. At the size we measure, one call of the translated C version takes between
39 and 789\,ms, with a median of 100\,ms.

\section{Transformation Classes}
\label{sec:loop-kernel-classes}

\Cref{tab:loop-classes} groups the kernels by the transformation they test. The benchmark
declares the class of 31 kernels; we derived it for the other nine. A manual check agrees with
six of our nine labels and questions three. Two of the declared labels are doubtful as well:
\texttt{vpvts} and \texttt{vtvtv} are filed under control flow but contain no branch. We keep
all labels as they are, so that our results stay comparable with the benchmark's.

\begin{table}[t]
  \centering
  \small
  \begin{tabular}{@{}lrp{7.6cm}@{}}
    \toprule
    Class & $n$ & Kernels \\
    \midrule
    Reduction            & 10 & s311, s3110, s3111, s3112, s316, s318, s319, argmax\_with\_index, quasi\_affine\_reduce\_odd, scan\_affine\_decay \\
    Loop interchange     &  5 & s1232, s2233, s231, s233, s235 \\
    Control flow         &  5 & s2710, s275, vpvts, vtvtv, ext\_break\_capture \\
    Indirect addressing  &  4 & s4112, vag, scatter\_accum\_dup, segment\_reduce\_ragged \\
    Loop fusion          &  3 & fuse\_diamond, fuse\_move\_ifs, fuse\_stencil\_through\_transient \\
    Dependence distance  &  2 & ext\_war\_unit, versioned\_distance\_update \\
    Linear dependence    &  2 & s115, s119 \\
    Scalar expansion     &  2 & s252, s255 \\
    Wavefront            &  2 & wf\_diff\_skew, wf\_triangular \\
    Loop distribution    &  1 & s2275 \\
    Node splitting       &  1 & s1244 \\
    Interprocedural data flow & 1 & s152 \\
    Recurrence           &  1 & s323 \\
    Packing              &  1 & compact\_threshold\_pack \\
    \bottomrule
  \end{tabular}
  \caption{The 40 loop-transformation kernels by class. TSVC kernels are named by their TSVC
    identifier (\texttt{s311} is \texttt{tsvc\_2\_s311} in the benchmark).}
  \label{tab:loop-classes}
\end{table}

The classes are uneven: reduction has ten kernels, but nine of the fourteen classes have at most
two. We therefore report results per kernel and use the classes only to group them.
